\section*{Hoofdstuk \ref{ch:b3:climate-models} --- Klimaatmodellen}

\begin{solution}{exo:b3:climate-models:1}
$A = 14.0 - 1.5 = \num{12.5}$ en $B = 0.14 - 0.05 = \num{0.09}$, in
\unit{kcal.cm^{-2}} per maand. Een maand is \qty{2.63e6}{s}, dus
\qty{1}{kcal.cm^{-2}} per maand is $\num{4.184e7}/\num{2.63e6} =
\qty{15.91}{W/m^2}$: $A = \qty{199}{W/m^2}$, $B =
\qty{1.43}{W.m^{-2}.\celsius^{-1}}$. Bij \qty{15}{\celsius}, $199 + 21.5 =
\qty{220}{W/m^2}$.
\end{solution}

\begin{solution}{exo:b3:climate-models:2}
$(238.0 - 211.9)/1.816 = \qty{14.4}{\celsius}$; met 0.31, $(234.6 -
211.9)/1.816 = \qty{12.5}{\celsius}$. Elke 0.01 albedo kost
\qty{3.4}{W/m^2} en $3.4/1.816 = \qty{1.9}{\celsius}$.
\end{solution}

\begin{solution}{exo:b3:climate-models:3}
$C = 1025\times3992\times60 = \qty{2.46e8}{J.m^{-2}.K^{-1}}$; gedeeld door
de \qty{3.156e7}{s} van een jaar, \qty{7.78}{W.yr.m^{-2}.K^{-1}}. Dan $\tau
= 7.78/1.816 = \num{4.3}\text{ jaar}$, en $7.78/1.33 = \num{5.8}\text{ jaar}$.
\end{solution}

\begin{solution}{exo:b3:climate-models:4}
ECS $= 3.93/1.33 = \qty{2.95}{\celsius}$; TCR $\approx 3.93/1.95 =
\qty{2.02}{\celsius}$; verhouding \num{0.68}.
\end{solution}

\begin{solution}{exo:b3:climate-models:5}
$Q = \qty{335}{W/m^2}$. Sneeuwbal: $(335\times0.38 - 211.9)/1.816 =
\qty{-46.6}{\celsius}$; warm: $(335\times0.70 - 211.9)/1.816 =
\qty{12.4}{\celsius}$, nog boven \qty{10}{\celsius}; ramp: $335(0.54 +
0.016\,T) = 211.9 + 1.816\,T$ geeft $T = \qty{8.7}{\celsius}$. $N' =
\qty{-1.816}{W.m^{-2}.\celsius^{-1}}$ bij de eerste twee, stabiel, en $-1.816
+ 335\times0.016 = \num{+3.54}$ bij de ramp, instabiel.
\end{solution}

\begin{solution}{exo:b3:climate-models:6}
De warme tak eindigt waar de temperatuur \qty{0}{\celsius} bereikt:
$0.70\,Q = A$, $Q = \qty{302.7}{W/m^2}$, een zonneconstante van
\qty{1211}{W/m^2}, \qty{11.0}{\percent} lager dan vandaag. De sneeuwbal smelt
waar $0.38\,Q = A - 20B$, $Q = \qty{462.1}{W/m^2}$, $S_0 =
\qty{1848}{W/m^2}$. De drempel is $B(T^* - T_F)/(Q_0(1 - \alpha_F))$,
het verschil tussen de huidige temperatuur $T^*$ en de temperatuur waarbij het
ijs begint zich uit te breiden, omgerekend naar zonlicht: $1.816\times14.37/238 =
\qty{11.0}{\percent}$ hier, $1.816\times4.37/238 = \qty{3.3}{\percent}$ in
de tekst.
\end{solution}

\begin{solution}{exo:b3:climate-models:7}
(a)~\qty{-45.5}{\celsius}. (b)~$340\times0.38 + F = 211.9 - 18.2$: $F =
\qty{64.5}{W/m^2}$, 16 keer de forcering van een verdubbeling. (c)~Zolang de
forcering duurt, springt de planeet naar de warme tak, $(238.0 + 64.5 -
211.9)/1.816 = \qty{50}{\celsius}$; zodra die wordt weggenomen koelt de planeet
langs de warme tak tot \qty{14.4}{\celsius}, die bestaat bij de huidige Zon:
de planeet keert niet terug naar de sneeuwbal.
\end{solution}

\begin{solution}{exo:b3:climate-models:8}
$\Sigma = 1.80/7.0 + 0.70/90 = \qty{0.2649}{yr^{-1}}$, $\Pi = 1.10\times0.70/
(7.0\times90) = \qty{1.222e-3}{yr^{-2}}$. De wortels geven $\tau_f =
\num{3.84}\text{ jaar}$ en $\tau_s = \num{213}\text{ jaar}$; de benaderingen,
$7.0/1.80 = \num{3.89}\text{ jaar}$, en $90\times1.80/(1.10\times0.70) =
\num{210}\text{ jaar}$. $a_f = 0.597$ exact, $1.10/1.80 = 0.611$
ongeveer. ECS $= 3.93/1.10 = \qty{3.57}{\celsius}$; TCR $\approx
3.93/1.80 = \qty{2.18}{\celsius}$.
\end{solution}

\begin{solution}{exo:b3:climate-models:9}
Noemer $B + \beta = \num{5.555}$. Evenaar: $(340\times1.221\times0.70 -
211.9 + 3.74\times12.81)/5.555 = \qty{22.8}{\celsius}$; pool:
$(340\times0.506\times0.38 - 211.9 + 47.9)/5.555 = \qty{-17.7}{\celsius}$.
Het transport verwijdert $3.74\times(22.8 - 12.8) = \qty{37}{W/m^2}$ bij de
evenaar en brengt $3.74\times30.6 = \qty{114}{W/m^2}$ naar de pool.
\end{solution}

\begin{solution}{exo:b3:climate-models:10}
(a)~$(1 - x^2)P_2' = 3x - 3x^3$, waarvan de afgeleide is $3 - 9x^2 = -6P_2$.
(b)~De transportterm is $-\mathcal D T_2(-6P_2) = 6\mathcal DT_2P_2$.
De balans $Q(1 + s_2P_2)(1 - \alpha) - A - B(T_0 + T_2P_2) - 6\mathcal
DT_2P_2 = 0$ geldt voor alle $x$ als het constante deel en het $P_2$-deel
wegvallen, wat $T_0$ en $T_2$ geeft. (c)~$T(0) - T(1) = T_2(-\tfrac12 - 1)
= -1.5\,T_2$. Zonder transport $T_2 = 340\times(-0.476)\times0.70/1.816 =
\qty{-62.4}{\celsius}$: een verschil van \qty{94}{\celsius}. (d)~$T_2 =
-55.3/1.5 = \qty{-36.9}{\celsius}$, dus $B + 6\mathcal D = 113.3/36.9 =
3.07$ en $\mathcal D = \qty{0.21}{W.m^{-2}.\celsius^{-1}}$.
\end{solution}

\begin{solution}{exo:b3:climate-models:11}
(a)~$r = \tfrac12\Sigma\big(1 \pm \sqrt{1 - 4\Pi/\Sigma^2}\big)$; met
$\sqrt{1 - \epsilon} \approx 1 - \epsilon/2$, de wortels zijn $\Sigma -
\Pi/\Sigma \approx \Sigma$ en $\Pi/\Sigma$. (b)~Met $T_d = 0$, $C\,\mathrm
dT/\mathrm dt = kt - (\gamma - \lambda)T$, waarvan de oplossing vanuit rust is $T =
k(t - \tau_f)/(\gamma - \lambda) + k\tau_f\mathrm e^{-t/\tau_f}/(\gamma -
\lambda)$, $\tau_f = C/(\gamma - \lambda)$. Voor $t \gg \tau_f$, $T \approx
kt/(\gamma - \lambda)$, en bij de verdubbeling, $kt = F_{2\times}$: TCR
$\approx F_{2\times}/(\gamma - \lambda)$. (c)~$\Sigma = 1.95/8.1 +
0.62/110 = \qty{0.2464}{yr^{-1}}$, $\Pi = 1.33\times0.62/(8.1\times110) =
\qty{9.25e-4}{yr^{-2}}$; $\tau_f = \num{4.12}\text{ jaar}$ en $\tau_s =
\num{262}\text{ jaar}$ (ongeveer $1/\Sigma = \num{4.06}\text{ jaar}$ en
$\Sigma/\Pi = \num{266}\text{ jaar}$); $4\Pi/\Sigma^2 = 0.061$.
\end{solution}

\begin{solution}{exo:b3:climate-models:12}
\sloppy (a)~Gemiddelde \qty{3.50}{\celsius}, standaardafwijking \qty{0.81}{\celsius}.
(b)~ECS $= 1.88 + 1.86\,X$; correlatie \num{0.992}; residuele
standaardafwijking \qty{0.11}{\celsius}. (c)~$1.88 + 1.86\times0.7 =
\qty{3.18}{\celsius}$, met een spreiding $\sqrt{(1.86\times0.1)^2 + 0.11^2} =
\qty{0.22}{\celsius}$, ongeveer een kwart van de spreiding van het ensemble.
(d)~Geen fysisch mechanisme koppelt $X$ aan de terugkoppelingen; $X$ is gevonden door
het uitproberen van veel waarneembare grootheden, zodat de correlatie toeval kan zijn; een vertekening
die alle acht modellen delen; modellen uit dezelfde familie geteld als
onafhankelijk.\par
\end{solution}

\begin{solution}{pb:b3:climate-models:1}
\textbf{1.} $Q_0 = 1360/4 = \qty{340}{W/m^2}$: de planeet onderschept het zonlicht met een schijf
van oppervlakte $\pi R^2$ en verspreidt het over
een bol van oppervlakte $4\pi R^2$.

\textbf{2.} $340\times0.70 = \qty{238.0}{W/m^2}$; $211.9 + 1.816\times15 =
\qty{239.1}{W/m^2}$.

\textbf{3.} $(238.0 - 211.9)/1.816 = \qty{14.4}{\celsius}$.

\textbf{4.} $(129.2 - 211.9)/1.816 = \qty{-45.5}{\celsius}$, onder
\qty{-10}{\celsius}.

\textbf{5.} $340(0.54 + 0.016\,T) = 211.9 + 1.816\,T$: $T =
(183.6 - 211.9)/(1.816 - 5.44) = \qty{7.8}{\celsius}$.

\textbf{6.} $(238.0/\num{5.670e-8})^{1/4} = \qty{254.5}{K}$,
\qty{-18.6}{\celsius}: \qty{33}{\celsius} onder het warme evenwicht, het
broeikaseffect van de planeet, ingebouwd in $A$ en $B$.

\textbf{7.} $N' = -B = \qty{-1.816}{W.m^{-2}.\celsius^{-1}}$ bij de
sneeuwbal en het warme evenwicht, beide stabiel; $-1.816 + 5.44 =
\num{+3.62}$ bij het middelste, instabiele evenwicht.

\textbf{8.} $\tau = 8.1/3.62 = \num{2.2}\text{ jaar}$.

\textbf{9.} $8.1/1.816 = \num{4.5}\text{ jaar}$.

\textbf{10.} $3.93/1.816 = \qty{2.2}{\celsius}$, onder de beoordeelde beste
schatting van \qty{3}{\celsius} en onder het \emph{waarschijnlijke} bereik van 2.5 tot
\qty{4}{\celsius}: de $B$ van het model, gefit over breedtegraden, is niet
de klimaatterugkoppelingsparameter van het klimaat als geheel.

\textbf{11.} $\mathrm e^{-t/\tau} = 0.05$ bij $t = 3\tau =
\num{13}\text{ jaar}$. Het echte klimaat warmt ook de diepe oceaan op, waarvan de tijdschaal
eeuwen is.

\textbf{12.} Bij $T_F = \qty{10}{\celsius}$; $Q = (211.9 + 18.2)/0.70 =
\qty{328.7}{W/m^2}$.

\textbf{13.} $S_0 = \qty{1315}{W/m^2}$, \qty{3.3}{\percent} onder dat van vandaag.

\textbf{14.} Naar de sneeuwbal, $(328.7\times0.38 - 211.9)/1.816 =
\qty{-47.9}{\celsius}$.

\textbf{15.} $Q = (211.9 - 18.2)/0.38 = \qty{509.8}{W/m^2}$, $S_0 =
\qty{2039}{W/m^2}$, \qty{50}{\percent} boven dat van vandaag.

\textbf{16.} Langs de warme tak omlaag naar \qty{10}{\celsius} bij
\qty{1315}{W/m^2}, een daling tot \qty{-48}{\celsius}, dan weer omhoog langs de
sneeuwbaltak naar \qty{-45.5}{\celsius} bij \qty{1360}{W/m^2}: de planeet
blijft bevroren. Pas boven \qty{2039}{W/m^2} zou hij terugspringen naar de warme
tak.

\textbf{17.} $\bar\alpha = 0.30\times0.9738 + 0.62\times0.0262 =
\num{0.3084}$; $\bar T = (340\times0.6916 - 211.9)/1.816 =
\qty{12.80}{\celsius}$.

\textbf{18.} $T_c = (340\times0.5643\times0.50 - 211.9 + 3.739\times12.80)/
5.555 = \qty{-12.26}{\celsius}$, \qty{3.5}{\celsius} warmer dan de
reanalyse bij \ang{72}.

\textbf{19.} $5.555\times(-12.26) + 211.9\times3.059 = \num{580.1}$, met
$\beta/B = \num{2.059}$.

\textbf{20.} De noemers zijn 1.7053, 1.7473, 1.7565, 1.7488 en
1.7362, en $Q/Q_0 = \num{1.0005}$, \num{0.9764}, \num{0.9713},
\num{0.9756} en \num{0.9827}.

\textbf{21.} Nabij \ang{47.5}, bij \num{0.971}. De ijslijnen poolwaarts van het
minimum zijn stabiel, $\mathrm dQ/\mathrm dx_s > 0$; die equatorwaarts zijn
instabiel.

\textbf{22.} $Q/Q_0 = 1$ ook tussen \ang{30} en \ang{40}, dicht bij
\ang{30} (\ang{30.2} door interpolatie): een instabiele ijslijn. Bij \ang{0}
is $Q/Q_0 = 580.1/(0.6106 + 0.7824)/340 = \num{1.225}$: de sneeuwbal houdt
stand voor elke Zon tot \qty{22}{\percent} helderder dan die van vandaag, dus
hij is vandaag mogelijk.

\textbf{23.} \qty{2.9}{\percent}, tegen \qty{3.3}{\percent} voor het
nuldimensionale model en \qty{1.6}{\percent} voor Budyko; het ijs staat
bij de drempel nabij \ang{48}, tegenover zijn \ang{50}. Het breedtegraadmodel
vergletsjert eerder dan het globale, omdat het zonlicht het zwakst is waar
het ijs zich vormt.

\textbf{24.} De modelplaneet vergletsjert bij een daling van
\textbf{\qty[text-series-to-math=true]{2.9}{\percent}} van de
zonneconstante (\qty{2.87}{\percent} volgens de berekening).
\end{solution}
